\documentclass[a4paper,11pt]{article}

\usepackage[a4paper, margin=2.5cm]{geometry}
\usepackage{booktabs}
\usepackage{siunitx}
\usepackage{xcolor}
\usepackage{colortbl}
\usepackage{caption}
\usepackage{multirow}
\usepackage{amsmath}
\usepackage{microtype}
\usepackage{parskip}
\usepackage{hyperref}

\definecolor{rowgray}{gray}{0.93}
\definecolor{highlight}{RGB}{220,240,220}
\definecolor{warnred}{RGB}{255,220,220}

\sisetup{
  group-separator  = {,},
  group-minimum-digits = 4,
  round-mode       = places,
  round-precision  = 1,
}

\title{\textbf{KissFeatures Message Size Analysis}\\[0.4em]
       \large MRC Multi-Robot Calibration --- Proximity Test (2 robots, Setup~1)}
\author{MRC Message Size Comparison Script}
\date{ROS~2 Jazzy, ec\_swift robot, RGL lidar simulation}

\begin{document}
\maketitle

% ─────────────────────────────────────────────────────────────
\section{Background}
% ─────────────────────────────────────────────────────────────

The MRC (Multi-Robot Calibration) slave node subscribes to raw
\texttt{sensor\_msgs/PointCloud2} lidar topics, merges them, and on a service
call returns a \texttt{KissFeatures} message containing FasterPFH keypoints
and descriptors to the MRC master for registration.

This report quantifies how large the \texttt{KissFeatures} message is compared
to the raw lidar data it was derived from, and what size reduction was achieved
by optimising the message encoding.

\textbf{Test setup:}
2 ec\_swift robots placed side-by-side ($\approx$1--2\,m apart).
Each robot carries two RGL lidars (front + back), providing 360\textdegree{} coverage.
Measurements are taken from robot\,0 as representative values
(5 samples per topic, 5 service calls).

% ─────────────────────────────────────────────────────────────
\section{Raw PointCloud2 Messages}
% ─────────────────────────────────────────────────────────────

Table~\ref{tab:raw_pc2} shows the serialised size and point count of each raw
lidar topic published by robot\,0.

\begin{table}[h]
\centering
\caption{Raw \texttt{PointCloud2} messages -- robot\,0}
\label{tab:raw_pc2}
\begin{tabular}{lSSSS}
\toprule
\textbf{Topic} &
  {\textbf{Samples}} &
  {\textbf{Avg size (KB)}} &
  {\textbf{Avg points}} &
  {\textbf{Max points}} \\
\midrule
\rowcolor{rowgray}
\texttt{front\_lidar/points\_raw} & 5 & 141.9 & 9073 & 9116 \\
\texttt{back\_lidar/points\_raw}  & 5 & 151.0 & 9651 & 9704 \\
\midrule
\rowcolor{highlight}
\textbf{Combined (per snapshot)} & & \textbf{292.9} & \textbf{18724} & \\
\bottomrule
\end{tabular}
\end{table}

Each point occupies exactly \SI{16}{\byte}: 12\,B for $xyz$ (3\,$\times$\,\texttt{float32})
plus 4\,B intensity (\texttt{float32}).
This is the data the slave must process every accumulation cycle.

% ─────────────────────────────────────────────────────────────
\section{KissFeatures Message Breakdown}
% ─────────────────────────────────────────────────────────────

The slave extracts FasterPFH keypoints and descriptors from the merged lidar
cloud and returns them as a \texttt{KissFeatures} message.
Table~\ref{tab:kiss_fields} breaks down each field in the message before and
after the encoding optimisation.

\begin{table}[h]
\centering
\caption{\texttt{KissFeatures} field sizes (avg over 5 service calls,
         $\approx$3387 keypoints, descriptor dim~=~33)}
\label{tab:kiss_fields}
\renewcommand{\arraystretch}{1.25}
\begin{tabular}{p{4cm} p{3.5cm} S[table-format=4.1] S[table-format=4.1] S[table-format=3.1]}
\toprule
\textbf{Field} &
  \textbf{Type} &
  {\textbf{Orig. (KB)}} &
  {\textbf{Opt. (KB)}} &
  {\textbf{Saving (\%)}} \\
\midrule
\rowcolor{rowgray}
\texttt{keypoints} (PointCloud2) &
  \texttt{sensor\_msgs/}\newline\texttt{PointCloud2} &
  53.0 & {---} & {---} \\

\texttt{keypoint\_xyz} (flat) &
  \texttt{float32[]} &
  {---} & 39.7 & {---} \\

\rowcolor{rowgray}
\texttt{keypoint\_count} &
  \texttt{uint32} &
  0.0 & {---} & {---} \\

\texttt{descriptors} &
  \texttt{float32[]} $\to$ \texttt{uint8[]} &
  437.0 & 109.0 & 75.1 \\

\rowcolor{rowgray}
\texttt{descriptor\_dim} &
  \texttt{uint32} &
  0.0 & 0.0 & {---} \\

\texttt{descriptor\_type} &
  \texttt{string} &
  0.0 & 0.0 & {---} \\

\texttt{header} &
  \texttt{std\_msgs/Header} &
  0.0 & 0.0 & {---} \\
\midrule
\rowcolor{highlight}
\textbf{Total} & &
  \textbf{490.0} & \textbf{148.9} & \textbf{69.6} \\
\bottomrule
\end{tabular}
\end{table}

\paragraph{Changes made.}
\begin{itemize}
  \item \texttt{sensor\_msgs/PointCloud2 keypoints} replaced by
        \texttt{float32[] keypoint\_xyz} (row-major $N \times 3$).
        Eliminates $\approx$13\,KB of PointCloud2 header and field-descriptor overhead.
  \item \texttt{uint32 keypoint\_count} removed (now derived as
        \texttt{keypoint\_xyz.size() / 3}).
  \item \texttt{float32[] descriptors} replaced by \texttt{uint16[] descriptors}.
        Encoding: $u = \mathrm{round}(v \times 655.35)$ (scale $= 65535/100$),
        decoding: $v = u / 655.35$.
        FasterPFH bins are normalised to sum to 100 per sub-histogram,
        so values lie in $[0, 100]$ and fit losslessly into 16\,bits
        (quantisation error $\leq 7.6\times10^{-4}$ per bin).
\end{itemize}

% ─────────────────────────────────────────────────────────────
\section{Message Size Comparison}
% ─────────────────────────────────────────────────────────────

Table~\ref{tab:comparison} compares the total wire size of the three data
representations for robot\,0's contribution to one calibration exchange.

\begin{table}[h]
\centering
\caption{Wire-size comparison for one calibration data exchange (robot\,0)}
\label{tab:comparison}
\begin{tabular}{lSSS}
\toprule
\textbf{Representation} &
  {\textbf{Size (KB)}} &
  {\textbf{vs.\ raw PointCloud2}} &
  {\textbf{vs.\ original KissFeatures}} \\
\midrule
\rowcolor{rowgray}
Raw PointCloud2 (2 lidars) & 292.9 & {1.0\,$\times$} & {0.6\,$\times$} \\
Original KissFeatures       & 490.0 & {1.7\,$\times$ \textbf{larger}} & {1.0\,$\times$} \\
\rowcolor{highlight}
\textbf{Optimised KissFeatures} & \textbf{148.9} &
  \textbf{2.0\,$\times$ smaller} & \textbf{3.3\,$\times$ smaller} \\
\bottomrule
\end{tabular}
\end{table}

% ─────────────────────────────────────────────────────────────
\section{Point-Count Reduction}
% ─────────────────────────────────────────────────────────────

Independently of byte size, the feature extraction step dramatically reduces
the number of geometric points that must be matched.

\begin{table}[h]
\centering
\caption{Point-count reduction through FasterPFH feature extraction}
\label{tab:points}
\begin{tabular}{lSS}
\toprule
\textbf{Stage} & {\textbf{Points}} & {\textbf{Bytes / point}} \\
\midrule
\rowcolor{rowgray}
Raw lidar (front + back combined) & 18724 & 16.0 \\
KissFeatures keypoints            &  3387 & {---} \\
\midrule
\rowcolor{highlight}
\textbf{Reduction} &
  \textbf{5.5\,$\times$} &
  {\textbf{81.9\,\% fewer}} \\
\bottomrule
\end{tabular}
\end{table}

The keypoint cloud is 5.5$\times$ sparser than the raw input, meaning the
registration algorithm (ROBIN matching + GNC-TLS) operates on roughly
\SI{18}{\percent} of the original points.

% ─────────────────────────────────────────────────────────────
\section{Summary}
% ─────────────────────────────────────────────────────────────

\begin{table}[h]
\centering
\caption{Summary of size reductions}
\label{tab:summary}
\begin{tabular}{lll}
\toprule
\textbf{Metric} & \textbf{Value} & \textbf{Notes} \\
\midrule
\rowcolor{rowgray}
Descriptor encoding      & float32 $\to$ uint8        & $\times$4 smaller per element \\
Descriptor array size    & 437\,KB $\to$ 109\,KB      & $-$75\,\% \\
\rowcolor{rowgray}
Keypoint representation  & PointCloud2 $\to$ float32{[]} & removes header overhead \\
Keypoint array size      & 53\,KB $\to$ 40\,KB        & $-$25\,\% \\
\rowcolor{rowgray}
Removed \texttt{keypoint\_count} & 4\,B $\to$ 0\,B    & now implicit in array size \\
\midrule
\rowcolor{highlight}
\textbf{Total KissFeatures} &
  \textbf{490\,KB $\to$ 149\,KB} &
  \textbf{$-$69.6\,\%, 3.3$\times$ smaller} \\
\rowcolor{highlight}
\textbf{vs.\ raw PointCloud2} &
  \textbf{1.7$\times$ larger $\to$ 2.0$\times$ smaller} &
  \textbf{flip from overhead to saving} \\
\bottomrule
\end{tabular}
\end{table}

The dominant cost before optimisation was the descriptor array (89\,\% of total
message size). Quantising the 33 FasterPFH bins from \texttt{float32} to
\texttt{uint16} reduces this by 50\,\% with negligible impact on matching
quality, since bin values lie in $[0,100]$ and the quantisation error
($\leq 7.6\times10^{-4}$ per bin) is negligible relative to the
inter-descriptor distances used by the ROBIN matcher.

After optimisation, transmitting \texttt{KissFeatures} is \textbf{2$\times$
cheaper} than sending the raw lidar scans, while providing a semantically
richer, descriptor-indexed representation that enables direct correspondence
search without any further processing by the master.

\end{document}
